(1) 令 \(m=|l|\ge1\)，则 \(|U|^2=|C|^2(r/w)^{2m}e^{-2r^2/w^2}\)。归一化
\[
1=\int_0^{2\pi}d\phi\int_0^\infty |U|^2r\,dr
=2\pi|C|^2{w^2m!\over2^{m+2}},
\]
故 \(|C|^2=2^{m+1}/(\pi w^2m!)\)。对强度求导 \(d\ln I/dr=2m/r-4r/w^2=0\)，得峰值环半径 \(r_\mathrm{peak}=w\sqrt{m/2}\)。
(2) 径向矩积分
\[
\langle r^{2s}\rangle={w^{2s}\over2^s}{(m+s)!\over m!}\quad(s=0,1,2,\dots).
\]
因此 \(\langle r^2\rangle=w^2(m+1)/2,\ \operatorname{Var}(r^2)=w^4(m+1)/4\)。当 \(m\gg1\)，峰与均方半径均按 \(w\sqrt{m/2}\) 增大，环的相对宽度缩小。
(3) 因 \(-i\hbar\partial_\phi U=l\hbar U\)，归一化后 \(\langle L_z\rangle=l\hbar\) 每个光子。变 \(l\to-l\) 时所有只依赖 \(m=|l|\) 的强度、径向矩不变，OAM方向反转。
(4) 当 \(m>0\)，\(U(0)=0\)，绕轴闭环相位变化 \(2\pi l\)，表现为相位涡旋。角动量密度与 \(l|U|^2\) 成正比，不意味着轴上有无限能量。有限孔径会混合径向模；高NA聚焦需矢量Maxwell场，SAM/OAM分解可能依赖选取的定义和参照系。
